% !TEX program = xelatex
\documentclass[UTF8, fontset=none, 11pt]{ctexart}
\setCJKmainfont{SimSun}
\setCJKsansfont{SimHei}
\usepackage[a4paper, margin=2.2cm]{geometry}
\usepackage{booktabs}
\usepackage{longtable}
\usepackage{array}
\usepackage{enumitem}
\usepackage{xcolor}
\usepackage{hyperref}
\hypersetup{colorlinks=true, linkcolor=black, urlcolor=blue!60!black}

\setlist{nosep, leftmargin=1.6em}
\renewcommand{\arraystretch}{1.3}

\title{\textbf{校园规则怪谈：自管委}\\[4pt]\Large 玩法详解暨交互手册（v2.1）}
\author{多人联网规则怪谈生存游戏}
\date{2026 年 11 月}

\begin{document}
\maketitle
\tableofcontents
\newpage

\section{游戏概述}
\subsection{世界观}
你加入了``学生自我管理委员会''（自管委）。委员会里混进了\textbf{不该存在的理事长}。干事们必须在一条条会被公开更换的校规下活下去，并把他们揪出来。

\subsection{设计支柱}
\begin{enumerate}
  \item \textbf{社区是玩法核心}：学生社区是推理、对抗与资源转化的主战场；其他空间 = 资源与重要物品探索、恐怖事件、位置推理。
  \item \textbf{违规记录双重属性}：干事的伤疤（SAN $-20$、被检举的证据）+ 理事长的弹药（当日检举额度）。理事长无数值无压力自由活动——``干净''本身引人怀疑，\textbf{理事长胜率偏高（55--60\%）是设计意图}。
  \item \textbf{懂规矩即战力}：通报隐去规则名，理事长必须熟记规则并用监控比对才能抓人；干事必须熟记规则，才能在理事长把全新规则换上板面后不踩新陷阱。
  \item \textbf{时间是资源}：全天 11 分钟四时段，复杂机制由客户端提示消化。
\end{enumerate}

\subsection{阵营与人数}
\begin{center}
\begin{tabular}{ll}
\toprule
项目 & 规则 \\
\midrule
玩家人数 & 9 起步（可扩 12）；潜伏理事长 3 名（9 人局基准） \\
\textbf{干事}（多数派） & 必须守规矩；HP/SAN/AP/饥饿四项数值；个人战 \\
\textbf{理事长}（少数派） & 无数值无压力；每日检举后可用备选池规则与板上 1 条随机规则相互替换（全校广播） \\
玩家资料 & 注册填\textbf{性别}（R1 判定需要） \\
身份 & 仅自己可见；淘汰公布阵营；死亡转``幽灵观察者'' \\
\bottomrule
\end{tabular}
\end{center}

\section{数值与资源}
\begin{center}
\begin{tabular}{llll}
\toprule
数值 & 初始 & 干事 & 理事长 \\
\midrule
生命值 HP & 100 & $\le 0$ 淘汰 & 无 \\
精神值 SAN & 100 & $\le 0$ 淘汰；$\le 60$ 轻度幻觉；$\le 30$ 重度幻觉 & 无 \\
行动值 AP & 10/天 & 移动 1、行动 2--4、睡觉 2、查阅 1 & 无 \\
饥饿值 & 100 & 每天 $-33$；$<40$ 扣 SAN、$<20$ 扣 HP、归 0 淘汰 & 无 \\
校园币 & 20 & 杂物回收、NPC 任务；私聊可转账 & 无 \\
\bottomrule
\end{tabular}
\end{center}

\textbf{幻觉}：SAN $\le 60$ 轻度（不存在的教室、四号窗口、错乱门牌）；SAN $\le 30$ 重度（投票噪声、文本污染）。幻觉由服务端下发，玩家不自知。

\section{通用违规后果（所有规矩适用）}
\begin{enumerate}
  \item SAN $-20$；生成违规记录（理事长弹药 $+1$）；
  \item \textbf{匿名通报}：校园墙广播``【通报】HH:MM，有学生违规。''——\textbf{只显示有人违规，不显示哪一条}；
  \item 部分规矩附带\textbf{即时惩罚}（当场结算）。
\end{enumerate}
\textbf{理事长的核心挑战}：通报只有时间、没有名字也没有规则名。检举\textbf{只需报出目标名字}，服务端核验其今日是否有违规记录——但``该抓谁''须自己推断：熟记 13 条规矩，用监控位置比对通报时间。

\section{一昼夜流程（11 分钟）}
\begin{center}
\begin{tabular}{llll}
\toprule
时段 & 真实时长 & 虚拟时刻 & 内容 \\
\midrule
上午 & 3 min & 08:00--11:00 & 探索（禁止在宿舍，F1） \\
中午 & 3 min & 11:00--13:00 & 自动传送食堂：讨论 + 投票 + 可选进餐 \\
下午 & 3 min & 13:00--17:00 & 探索 \\
晚上 & 2 min & 17:00--21:00 & 收尾；21:00 前回宿舍（F5） \\
深夜 & --- & 21:00--23:00 & 理事长内部交流 + F7 呼唤 \\
\bottomrule
\end{tabular}
\end{center}
AP 每天 10；前往食堂自动传送不耗 AP。\textbf{Day 0}：中午不投票、无检举、无更换、社区自由进出。

\subsection{清晨}
查看昨夜结果与规则板公告 $\to$ \textbf{社区预约}（上午/下午场各限 3 人，每人每天限 1 场；Day 1 起未预约进入 $=$ 违规``翻越闸机''）$\to$ 理事长弹药刷新（$=$ 当日违规记录数，共享池）。

\subsection{上午 / 下午：探索 + 检举}
移动 1 AP；A 类节点内移动 1 AP（教学楼按楼层移动：上一层楼 1 AP）；各子位置探索 2--4 AP（文案见附录 A）。\textbf{检举}见第 9 节。

\subsection{中午：讨论 + 投票（3 分钟）}
自动传送食堂 $\to$ 可选进餐（$\sim$3 币 $+35$ 饥饿，可能``食物有问题''）$\to$ 自由发言（表彰者定发言顺序）$\to$ 匿名投票（表彰者 1.5 票 + 归票权）$\to$ 最高票出局公布阵营，平票平安夜。

\subsection{晚上与深夜}
21:00 未回宿舍 $=$ F5 违规。深夜理事长内部频道交流；F7 呼唤事件随机降临（回应 $=$ 违规）。

\section{地图（14 节点）与可见性}
\textbf{节点}：宿舍、食堂、教学楼 A/B/C、超市、图书馆、校医院、操场、学生社区、回收站、行政楼、学校大门、体育馆。

\textbf{连接（1 AP）}：宿舍--食堂、宿舍--学校大门、宿舍--学生社区、宿舍--校医院、校医院--食堂、食堂--教学楼 A、食堂--超市、食堂--操场、食堂--学生社区、教学楼 A--B、B--C、A--图书馆、图书馆--行政楼、行政楼--学校大门、超市--回收站、回收站--体育馆、操场--体育馆。

\textbf{A 类子位置}（进入落在大厅/前台，节点内移动 1 AP）：\textbf{教学楼 A（5 层 17 间）、B（4 层 15 间）、C（5 层 15 间）}——教室编号 $=$ 楼层+序号（如 301），每间独立文案（附录 A.3--A.5）；图书馆（大厅、阅览室、书库区、自习区）；行政楼（大厅、影像室、网络中心、文印室）；体育馆（前台、器械区、球场、储物间）；学生社区（前台、办公室、社区部、策创部、宣传部、咖啡屋）。B 类（宿舍、食堂、操场、超市、校医院、回收站、学校大门）无子位置但有多个可探索点。

\textbf{可见性}：同节点玩家互相可见到具体房间；跨节点不可见；断网期间全局不可见。目击记录按房间精度生成。

\section{学生社区（预约制）}
钥匙持久可交易、每种每局限量 2 把；房间每人每天限 1 次；\textbf{进哪个房间同社区者可见}。
\begin{center}
\begin{tabular}{lll}
\toprule
房间 & 钥匙 & 功能 \\
\midrule
办公室 & 办公室钥匙 & 档案 $\times 1$：指定 2 人验是否同阵营；每日全图仅 1 份 \\
社区部 & 社区部钥匙 & 昨日社区预约与出入时间线 \\
策创部 & 策创部钥匙 & 抹除 1 条随机规则（每人每局限 1 次） \\
宣传部 & 宣传部钥匙 & 校园墙匿名发帖 1 条（破译 / 表彰 / 担保） \\
咖啡屋 & 咖啡券 & 50 HP 或 30 SAN（二选一）+5 饥饿 \\
\bottomrule
\end{tabular}
\end{center}

\textbf{校园墙三件套}：（1）\textbf{破译}：选 2 条规矩永久详细化；（2）\textbf{三好学生表彰}：被表彰者当天定发言顺序、1.5 票、归票权，干事零违规则次日 AP 15（理事长同样生效，无 AP 加成）；（3）\textbf{匿名担保}：被担保人今日已违规则当晚免疫检举，否则\textbf{双方各 $-25$ HP / $-25$ SAN}。

\section{行政楼（信息中枢）}
\begin{center}
\begin{tabular}{ll}
\toprule
房间 & 功能 \\
\midrule
影像室 & 干事每日 3 次查全图节点级位置；理事长不限；\textbf{双方留痕} \\
网络中心 & 剪网线（剪刀 1.5 min）/ 杠铃杆破坏（4 min）/ 修复（1 min 或半天自恢复） \\
文印室 & 查阅卡产出（每日 1 张，先到先得，持久可交易） \\
\bottomrule
\end{tabular}
\end{center}

\textbf{机器查阅}：持查阅卡查当日记录——超市售货机、图书馆/操场/社区闸机、影像室调取。每次 1 AP + SAN 3，\textbf{只有时间次数无人名}。\textbf{留痕推理法}：影像室调取 $> 3 \times$ 存活干事数 $\Rightarrow$ 超出必属理事长。

\section{校园规矩体系}
规则板 \textbf{13 条 $=$ 8 固定 + 5 随机}（12 池抽取）；每条规矩\textbf{必须可通过操作或配合避免}。

\subsection{固定规矩（8 条）}
\begin{longtable}{p{0.7cm}p{2.6cm}p{5.8cm}p{5.9cm}}
\toprule
编号 & 名称 & 规则板文本 & 判定与即时惩罚 \\
\midrule
\endhead
F1 & 上午宿舍清空 & 【8:00--11:00】上午，禁止出现在宿舍。 & 上午结束位置 $=$ 宿舍 $\to$ 违规。（睡觉限 11:00 后） \\
F2 & 禁止打扰教学 & 特殊期间无需上课，禁止打扰老师的教学活动。 & 进入「上课中」教室 $\to$ 违规；\textbf{即时：逐出教学楼 + SAN $-10$}。 \\
F3 & 教学楼禁食 & 教学楼内禁止进食。 & 位于教学楼 A/B/C 内「进食」$\to$ 违规。 \\
F4 & 教学区禁打闹 & 教学楼与图书馆内禁止打闹。 & 在两处使用杠铃杆 $\to$ 违规；\textbf{即时：没收 + HP $-10$}。 \\
F5 & 禁止夜不归宿 & 21:00 前必须回到宿舍。恰好归零时回到也算合规。 & 21:00 位置 $\ne$ 宿舍 $\to$ 违规；AP 归零不在宿舍 $\to$ 力竭滞留（去重）。 \\
F6 & 禁止进入不存在的教室 & 教室编号为``楼层+序号''（如 301）。\textbf{404 不存在。} & SAN $<60$ 幻觉门（404、$-1$ 层、13 号）进入 $\to$ 违规 + SAN $-15$。 \\
F7 & 熄灯后不应答呼唤 & 21:00 后听到有人叫你的名字，不要回应。 & 深夜回应「呼唤」$\to$ 违规（物资诱惑）。 \\
F8 & 禁止空票 & 每天中午的投票必须投出一票。 & 投票结束未提交 $\to$ 违规。 \\
\bottomrule
\end{longtable}

\subsection{随机规矩池（12 选 5）}
\begin{longtable}{p{0.7cm}p{2.4cm}p{5cm}p{6.2cm}}
\toprule
编号 & 名称 & 规则板文本 & 判定 \\
\midrule
\endhead
R1 & 男女不得独处 & 男女生不得单独相处。 & A 类节点\textbf{同房间仅剩两人且异性} $\to$ 违规（B 类不适用；可拉第三人进房或离开）。 \\
R2 & 禁止尾随 & 不要跟着别人。 & 同日第二次与同一玩家先后进同一房间（$\pm3$ min）$\to$ 违规。 \\
R3 & 体育馆宵禁 & 19:00 后体育馆关闭。 & 19:00 后位于体育馆 $\to$ 违规。 \\
R4 & 图书馆午休 & 图书馆 11:00--13:00 闭馆。 & 该时段位于图书馆 $\to$ 违规。 \\
R5 & 宿舍禁食 & 宿舍内禁止进食。 & 位于宿舍时「进食」$\to$ 违规。 \\
R6 & 校门封闭 & 白天禁止在学校大门逗留。 & 白天位于学校大门 $\to$ 违规。 \\
R7 & 行政楼准入 & 生命值高于 80 者禁止进入行政楼。 & 进入时 HP $>80$ $\to$ 违规。 \\
R8 & 限时安静 & 上午和下午公共频道发言各不得超过 8 条。 & 服务端分时段计数。 \\
R9 & 红校服礼节 & 见到红色校服的同学，点头致意即可，不要搭话。 & NPC 事件选「搭话」$\to$ 违规。 \\
R10 & 讨论禁止离场 & 午间讨论期间不要离开食堂。 & 11:00--13:00 离开食堂节点 $\to$ 违规。 \\
R11 & 精神自检 & SAN 低于 40 的同学每天必须用药或运动。 & 21:00 时 SAN $<40$ 且当日未用药/运动 $\to$ 违规。 \\
R12 & 门牌确认 & 每天回到宿舍后，请确认门牌号。 & 确认事件：SAN $<60$ 时门牌错误，应点「报修」。 \\
\bottomrule
\end{longtable}

\subsection{规则板运作（更换制 $=$ 整规则互换）}
\begin{enumerate}
  \item \textbf{检举操作员}：理事长阵营的检举由\textbf{编号最小的理事长}统一操作。
  \item \textbf{更换权}：每日检举结算后，操作员可将板上 1 条随机规则\textbf{与备选池中的 1 条规则相互替换}——只能从 12 池中\textbf{当前不在板上}的规则挑选（开局 5 条在板、7 条备选），\textbf{每日至多 1 次，也可放弃}。
  \item \textbf{更换即广播}：``【规则修正】即日起《R2 禁止尾随》废止，替换为：《R12 门牌确认》——每天回到宿舍后，请确认门牌号。''——新旧规则名与全文全校公示，无秘密篡改；换下的规则回到备选池，可被日后再次换上。
  \item \textbf{校规便签}（原校规原文卡）：揭示板上 1 条随机规矩的\textbf{隐藏判定细则}（一句话规则背后的完整判定说明，单条弱化版破译）。策创部仍可抹除 1 条。
  \item 每次更换都\textbf{引入一条全新规则}，干事须即时重读新规则；理事长可借此把``最容易诱捕''的规则换上板面。
\end{enumerate}

\section{检举弹药制}
\begin{itemize}
  \item \textbf{额度}：当日每 1 条学生违规记录 $=$ 阵营共享 1 次检举机会（客户端实时显示）。
  \item \textbf{判定}：操作员\textbf{只需报目标名字}（无需指认规则），服务端核验该干事今日是否存在违规记录。
  \item \textbf{伤害}：$25 + 5 \times day$（day 0 起算）。
  \item \textbf{停止}：操作员失败累计 2 次当日停止；额度用尽即止；无追加机制。
  \item 担保/护身符可免疫一次；理事长只见记录数不见人，通报又隐去规则名——\textbf{熟记规矩 + 监控比对}是唯一抓手。
\end{itemize}
\textbf{闭环}：干事守规矩 $\to$ 理事长没弹药 $\to$ 理事长换上诱捕规则制造陷阱 $\to$ 干事重读新规、策创部抹除。

\section{理事长特权}
\begin{enumerate}
  \item \textbf{假装交互}：房间交互有界面无实效（不耗次数、无资源、不可拾取）；\textbf{进房可查看物品清单与价值}（看得见拿不到）。
  \item \textbf{影像室}不限次（留痕）；\textbf{超市}免费伪装采购（不能吃）；\textbf{免钥匙}进社区房间（不可用功能）。
  \item \textbf{能力}（每人 1 个，单理事长局 2 个）：窥探者 / 伪证者 / 审判官 / 幕后者。
\end{enumerate}

\section{破坏、攻击与信息渠道}
\textbf{剪网线}（剪刀 30 币，1.5 min）/ \textbf{杠铃杆破坏网络中心}（4 min）：掐断校园墙、私聊、理事长群聊、监控与全局视野；行政楼 1 min 修复或半天自恢复。\textbf{攻击}（杠铃杆，同房间）：眩晕 2 min，干事额外 $-30/-30$，被目击；教学楼/图书馆内使用 $=$ F4 违规没收；\textbf{理事长不能持械}。

\begin{center}
\begin{tabular}{ll}
\toprule
频道 & 说明 \\
\midrule
公共频道 / 私聊（转账） & 全员 / 1 对 1 \\
理事长内部群聊 & 仅理事长；可被断网切断 \\
校园墙 & 只读；宣传部可匿名发帖；预约名单、通报、规则修正公告 \\
影像室监控 / 机器查阅 & 节点级位置（干事 3 次/日）；各场地记录（时间+次数无人名） \\
系统广播 & 不可切断 \\
\bottomrule
\end{tabular}
\end{center}

\section{生存经济与胜负}
饥饿每天 $-33$（$<40$ 扣 SAN、$<20$ 扣 HP、归 0 淘汰）。食堂午餐 $\sim$3 币 $+35$（非强制）、超市零食 10 币、咖啡券 5 币、校医院急救包/镇定剂各 40 币、体育馆 10 币、操场免费、宿舍睡觉 2 AP（SAN+20，仅 11:00 后）。杂物 3--8 币/件；重要物品交 NPC 换 20--50 币、钥匙、咖啡券或推理情报。

\textbf{胜负}：（1）理事长全灭 $\to$ 干事阵营大获全胜；（2）第 X 天结算：理事长存活 $\ge$ 干事存活 $\to$ 理事长大获全胜；（3）否则存活干事个人获胜。

\clearpage
\appendix
\section{全节点交互文案设计}
\label{app:rooms}
每个可探索子位置给出：入口文本、交互选项与结果文本、消耗与产出。文案为示例底稿，正式上线可按美术风格润色。幻觉类条目仅在 SAN $<60$ 时进入对应玩家的选项池。

\subsection{宿舍（B 类）}
\begin{itemize}
  \item \textbf{床铺 · 睡觉}（2 AP，仅 11:00 后）：``你拉上窗帘。黑暗里，走廊的脚步声停了。'' $\to$ SAN $+20$，黑屏 2 小时，饥饿 $\times 2$。
  \item \textbf{书桌 · 搜索}（2 AP）：``抽屉深处卡着一张上届留下的课表，周三下午是空白。'' $\to$ 杂物 / 重要物品「空白课表」（交 NPC 换情报）。
  \item \textbf{门牌 · 确认}（每日首次回宿舍）：``你的门牌号：308。要确认吗？''（SAN $<60$ 时显示错误号）$\to$ 选「确认无误 / 报修」$\to$ R12 判定。
  \item \textbf{阳台}（幻觉事件，SAN $<60$）：``对面楼的窗帘后站着一个人，姿势和你一样。'' 选项：「挥手」（SAN 判定，败 $-15$）/「拉上窗帘」（SAN $+5$）。
  \item \textbf{室友的空床}（事件）：``床上有起伏，像有人躺在被子里。'' 选项：「掀开」（SAN 判定：成 $+5$``只是衣服''；败 $-15$）/「无视」。
\end{itemize}

\subsection{食堂（B 类）}
\begin{itemize}
  \item \textbf{打饭窗口 · 午餐}（$\sim$3 币）：``今日特餐：红烧狮子头。吃吗？'' $\to$ $+35$ 饥饿；25\% 追加``饭菜里有个硬物，像是指甲。''（$-5\sim-15$ HP/SAN）。
  \item \textbf{餐桌区}（事件）：``一个不存在的同学端着托盘：这里有人吗？'' 「有人」（SAN $+5$）/「没有人」（SAN 判定，大失败 $-20$）/「低头吃饭」（SAN $+5$）。
  \item \textbf{后厨门}（事件）：``员工通道，门没锁，里面有汤的香气。'' 「推开」（SAN 判定：成 杂物 $+5$ 币；败 $-15$）/「离开」。
  \item \textbf{四号窗口}（幻觉，SAN $<60$）：``4 号窗口：回头客免费。'' 选择 $\to$ R9 违规 + SAN $-10$。
  \item \textbf{公告黑板}（免费）：``今日值日表。'' 你的名字被写在昨天的格子里。 $\to$ 情报碎片（纯叙事，可交 NPC）。
\end{itemize}

\subsection{教学楼 A（5 层 17 间，标准型）}
大厅公告栏显示各教室实时状态（空闲 / \textbf{上课中}）——F2 的官方信息源。上一层楼 1 AP。除标注外，搜索均为 3 AP。
\begin{longtable}{p{1.1cm}p{1.2cm}p{6.6cm}p{5.5cm}}
\toprule
房间 & 类别 & 入口文本 & 交互与产出 \\
\midrule
\endhead
101 & 物资 & ``教室整洁得反常，每张课桌左上角都摆着一块橡皮，摆得一样齐。'' & 杂物 $\times 2$。 \\
102 & 上课 & 门牌「上课中」。门上的课程表写着《观察者素养》。里面有老师在讲课。 & 进入 $\to$ F2：逐出 + SAN $-10$。 \\
103 & 物资 & ``讲台抽屉里有一沓饭票，日期全是明天。'' & 「明日饭票」$\times 1$（次日午餐免费）。 \\
201 & 上课 & 「上课中」。从门缝看，教室坐满了；推开门，一个人也没有。 & 进入 $\to$ F2。 \\
202 & 物资 & ``图书角的书全部倒着插，像有人从里面往外看。'' & 笔记 $\times 1$。 \\
203 & 挑战 & 黑板上有一道题：``此刻，这间教室里有几个人？'' & 答``1''（你自己）$\to$ 咖啡券 $\times 1$；错 SAN $-10$。 \\
204 & 物资 & ``窗台上半瓶汽水，没开封，还是凉的。'' & 零食（$+20$ 饥饿）。 \\
301 & 物资 & ``储物柜里挂着一件校服，尺码标签写着你的名字。'' & SAN 判定：败 SAN $-10$；成 杂物。 \\
302 & 上课 & 「上课中」。老师在点名，每点到一个名字，教室里就多一声应答。 & 进入 $\to$ F2。 \\
303 & 挑战 & 讲台上有一份问卷：``你今天违反了几条规矩？'' & 如实填写；填 0 且当日确为 0 $\to$ \textbf{护身符} $\times 1$；否则 SAN $-10$。 \\
401 & 物资 & ``404 的门牌挂在 401 旁边——牌子后面是一堵墙。''（SAN $<60$ 时牌子变成一扇门 $\to$ F6） & 笔记 $\times 1$。 \\
402 & 幻觉 & 仅 SAN $<60$ 可见此门。门内是一间坐满``你''的教室，所有``你''同时回头。 & 进入 $\to$ SAN $-15$；离开 $\to$ SAN $+5$。 \\
403 & 上课 & 「上课中」。这节课讲的内容，是你入学以来的档案。 & 进入 $\to$ F2。 \\
501 & 挑战 & 黑板上写着：``$\int_0^1 x\,dx = \,?$'' & 答 $1/2$ $\to$ 10 币；错 SAN $-10$。 \\
502 & 物资 & ``走廊尽头的天台门没锁——虽然这是五楼的最里面。'' & 重要物品「旧照片」。 \\
503 & 上课 & 「上课中」。下课铃从你推门那一刻开始响，一直没停。 & 进入 $\to$ F2。 \\
\bottomrule
\end{longtable}

\subsection{教学楼 B（4 层 15 间 + 幻觉位，实验艺术型）}
\begin{longtable}{p{1.1cm}p{1.2cm}p{6.6cm}p{5.5cm}}
\toprule
房间 & 类别 & 入口文本 & 交互与产出 \\
\midrule
\endhead
101 & 挑战 & 烧杯里泡着一张纸条：``什么东西越分享越多？'' & 答``知识''$\to$ 5 币；错 SAN $-10$。 \\
102 & 物资 & ``实验器材按高矮排成一列，最高的那只烧杯在发抖。'' & 杂物 $\times 2$。 \\
103 & 上课 & 「上课中」。 & 进入 $\to$ F2。 \\
104 & 挑战 & ``1+1 在什么情况下等于 3？'' & 答``算错的时候''$\to$ 咖啡券；错 SAN $-10$。 \\
201 & 幻觉 & 地板变成水面，天花板是另一间教室的倒影。（SAN $<60$） & 进入 $\to$ SAN $-10$；退出 $\to$ SAN $+5$。 \\
202 & 物资 & ``模型骨架朝你挥手——大概是风。'' & 杂物；低概率镇定剂。 \\
203 & 挑战 & 局势题：``昨天中午的投票，出现平票了吗？''（以系统记录为准） & 答对 $\to$ 10 币；错 SAN $-10$。 \\
204 & 上课 & 「上课中」。 & 进入 $\to$ F2。 \\
301 & 挑战 & 规则题：``规则板上现在有几条随机规矩？'' & 答``5''（按当前板面）$\to$ 校规便签 $\times 1$；错 SAN $-10$。 \\
302 & 物资 & ``画室里的肖像都没有脸，除了最新的一幅——那幅在画你。'' & SAN 判定：败 $-10$；成 杂物。 \\
303 & 上课 & 「上课中」。 & 进入 $\to$ F2。 \\
304 & 物资 & ``纸箱上写着：历届学生会遗落物。'' & 重要物品概率提升。 \\
401 & 挑战 & ``什么东西打破了才能吃？'' & 答``鸡蛋''$\to$ 零食 $\times 2$；错 SAN $-10$。 \\
402 & 幻觉 & 走廊尽头出现写着「404」的门。这层楼只有 401--403。（SAN $<60$） & 推门 $\to$ F6 违规 + SAN $-15$。 \\
403 & 上课 & 「上课中」。 & 进入 $\to$ F2。 \\
电梯 & 幻觉 & ``面板上多出一个「0 层」按钮。''（SAN $<60$，B 座通用） & 「按下」$\to$ HP $-10$ + SAN $-20$；「走楼梯」$\to$ SAN $+5$。 \\
\bottomrule
\end{longtable}

\subsection{教学楼 C（5 层 15 间，高风险高回报）}
搜索 4 AP，事件概率与重要物品概率提升。
\begin{longtable}{p{1.1cm}p{1.2cm}p{6.6cm}p{5.5cm}}
\toprule
房间 & 类别 & 入口文本 & 交互与产出 \\
\midrule
\endhead
101 & 物资 & ``桌上的考卷写满了正确答案，姓名栏是你的名字。'' & 笔记 $\times 1$。 \\
102 & 上课 & 「上课中」。 & 进入 $\to$ F2。 \\
103 & 挑战 & 问卷：``你认为理事会里有几个人？（答对有大奖）'' & 答 $=$ 实际理事长存活数 $\to$ 10 币 + SAN $-5$（``它知道你在猜''）；错 SAN $-10$。 \\
201 & 幻觉 & 教室里的钟表全部倒着走，粉笔自己在黑板上写字。（SAN $<60$） & 进入 $\to$ SAN $-10$。 \\
202 & 物资 & ``讲台下的暗格里有一枚铜钥匙，牌子上刻着`办公室'。'' & 50\% 真·办公室钥匙 / 50\% 玩具钥匙（SAN $-5$）。 \\
203 & 上课 & 「上课中」。点名声里，包含已经被淘汰的人。 & 进入 $\to$ F2。 \\
301 & 挑战 & 黑板上写着：``$e^{i\pi}+1=?$'' & 答 $0$ $\to$ \textbf{查阅卡} $\times 1$；错 SAN $-10$。 \\
302 & 物资 & ``墙皮剥落处刻着一行字：别在 C 座逗留太久。'' & 情报碎片。 \\
303 & 上课 & 「上课中」。 & 进入 $\to$ F2。 \\
401 & 幻觉 & 整间教室是你宿舍的样子，床铺铺好了，在等你。（SAN $<60$） & SAN 判定：败 SAN $-15$；成 SAN $+5$。 \\
402 & 挑战 & ``最后一题：你现在应该离开这栋楼吗？（是 / 否）'' & 无论选什么 $\to$ SAN $-5$ 并触发楼梯间事件。 \\
403 & 上课 & 「上课中」。 & 进入 $\to$ F2。 \\
501 & 物资 & ``天台的晾衣绳上挂着一排名牌，都是今年的新生。'' & 重要物品「名牌」（交 NPC 换情报）。 \\
502 & 挑战 & 局势题：``今天至今为止，校园墙通报过几次违规？'' & 答对 $\to$ 10 币；错 SAN $-10$。 \\
503 & 上课 & 「上课中」。黑板上写着：``下课后来办公室。'' & 进入 $\to$ F2。 \\
\bottomrule
\end{longtable}

\subsection{图书馆（A 类：大厅 / 阅览室 / 书库区 / 自习区）}
\begin{itemize}
  \item \textbf{阅览室 · 搜索}（3 AP）：``安静得能听见翻页声。'' $\to$ 笔记（高）/ 杂物；进出记入闸机记录。
  \item \textbf{书库区}（事件）：``有一格书脊上写着你的名字。'' 「抽出」（SAN 判定：成 笔记 $\times 1$；败 SAN $-20$）/「离开」。
  \item \textbf{自习区}（事件）：``邻座推来一张纸条：别相信规则。'' 「回信」（SAN 判定：成 NPC 情报；败 $-10$）/「不理」（SAN $+5$）。
  \item \textbf{闭馆广播}（事件）：``请以下同学立即离馆：……你的名字。'' 「离开」（放弃本次搜索）/「留下」（SAN 判定：成 免费追加一次搜索；败 SAN $-15$ 并被送出）。
\end{itemize}

\subsection{校医院（B 类）}
\begin{itemize}
  \item \textbf{药房 · 购药}：``处方还是零售？'' 急救包 40 币 / 镇定剂 40 币。
  \item \textbf{诊室 · 校医}（事件）：``喝了吧，对你有好处。'' 「喝」（HP $+10$ + SAN 判定：败``药里有东西在动'' SAN $-15$）/「拒绝」。
  \item \textbf{观察室}（事件）：``空病床上摆着一张体检单，照片是你，名字不是你。'' SAN 判定：败 SAN $-15$。
  \item \textbf{校医 NPC}：交付重要物品 $\to$ 20--50 币 / 咖啡券 / 情报。
\end{itemize}

\subsection{操场（B 类）}
\begin{itemize}
  \item \textbf{跑道 · 运动}（免费）：``风很干净。你跑了很久。'' $\to$ HP $+5$、SAN $+8$，消耗饥饿；进出记入闸机记录。
  \item \textbf{看台}（事件）：``夜跑的人数到第 13 圈，停下来看你。'' 「对视」（SAN 判定：败 $-15$）/「移开视线」（SAN $+5$）。
  \item \textbf{器材角 · 搜索}（2 AP）：$\to$ 杂物 / 体育物资（低）。
  \item \textbf{红校服}（R9 NPC，随机）：``红校服的同学朝你点头。'' 「点头」/「搭话」（违规）/「无视」（SAN $+5$）。
\end{itemize}

\subsection{体育馆（A 类：前台 / 器械区 / 球场 / 储物间）}
\begin{itemize}
  \item \textbf{前台 · 购票}（10 币）：``单次 10 币，请扫码。'' $\to$ 解锁器械区与球场。
  \item \textbf{器械区 · 运动}：``比操场更狠的训练。'' $\to$ HP $+12$、SAN $+10$，概率得体育物资（携带时饥饿加速下降）。
  \item \textbf{器械区 · 杠铃架}（3 AP）：``杠铃杆靠在墙角，握柄上有新鲜的指印。'' $\to$ \textbf{杠铃杆}（每人限持 1）。
  \item \textbf{球场}（事件）：``灯自己亮着，篮球在拍自己。'' 「捡起」（SAN 判定：成 体育物资 $\times 2$；败 SAN $-15$）/「离开」。
  \item \textbf{储物间 · 探索}（3 AP）：更衣柜 1--20，有一格没锁。 $\to$ 体育物资 / 杂物（回收价值高）。
  \item \textbf{13 号柜}（幻觉事件）：``你的柜门编号变成了 13。'' 「打开」（SAN $-15$）/「换柜」（SAN $+5$）。
\end{itemize}

\subsection{超市（B 类）}
\begin{itemize}
  \item \textbf{货架 · 购买}：``欢迎光临。'' 零食 10 币 / 剪刀 30 币 / 咖啡券 5 币。
  \item \textbf{自动售货机}：``哐当。'' 零食 $+15\sim25$ 饥饿；\textbf{每笔购买记入机器日志}（可被查阅卡查询）。
  \item \textbf{收银台 · 老板娘}（NPC）：``你看到了吗？'' 交付重要物品 $\to$ 20--50 币 / 钥匙 / 情报。
  \item \textbf{新到的货}（事件）：``生产日期是明天。'' 「购买」（SAN 判定：败 SAN $-5$）/「放下」。
\end{itemize}

\subsection{回收站（B 类）}
\begin{itemize}
  \item \textbf{秤台 · 卖杂物}（1 AP）：``大爷接过你的袋子，掂了掂。'' $\to$ 3--8 币/件。
  \item \textbf{废品堆 · 搜索}（2 AP）：$\to$ 杂物。
  \item \textbf{压扁的玩偶}（事件）：``玩偶怀里抱着字条：别去办公室。'' SAN 判定：成 情报线索；败 SAN $-10$。
  \item \textbf{回收站大爷}（NPC）：交付重要物品 $\to$ 报酬 / 情报。
\end{itemize}

\subsection{行政楼（A 类：大厅 / 影像室 / 网络中心 / 文印室）}
\begin{itemize}
  \item \textbf{大厅 · 公示屏}（免费）：``今日监控调取：N 次。''——留痕推理的信息源。
  \item \textbf{影像室 · 查询}：``调取今日监控？（今日剩余 N 次）'' $\to$ 全图节点级位置；\textbf{留痕}（干事 3 次/日，理事长不限）。
  \item \textbf{网络中心}：``机柜的灯一格格灭着。'' 剪刀剪网线（1.5 min）/ 杠铃杆破坏（4 min）/ 修复（1 min 或半天自恢复）。
  \item \textbf{文印室 · 领卡}：``复印机上压着一叠空白卡。'' $\to$ 查阅卡 $\times 1$（每日 1 张，先到先得）。
  \item \textbf{复印机}（幻觉事件）：``吐出一张写满你名字的纸。'' 「拿走」（SAN 判定：成 查阅卡线索；败 SAN $-15$）/「不理会」。
  \item \textbf{R7}：进入时 HP $>80$ $\to$ 违规（闸机红灯亮起，门卫抬头看了你一眼）。
\end{itemize}

\subsection{学校大门（B 类）}
\begin{itemize}
  \item \textbf{门卫室 · 门卫}（NPC）：``今天没有请假单。'' 交付重要物品 $\to$ 报酬 / 情报；白天逗留 $\to$ R6 违规。
  \item \textbf{铁门}（事件）：``你握住门把手，它是湿的。'' 「推门」（SAN 判定：败 SAN $-15$）/「松手」。
  \item \textbf{传达室窗}（免费）：``值班表上，今晚的值日生被划掉了三个名字。'' $\to$ 情报碎片。
\end{itemize}

\subsection{学生社区（A 类：前台 / 办公室 / 社区部 / 策创部 / 宣传部 / 咖啡屋 / 走廊）}
\begin{itemize}
  \item \textbf{前台 · 闸机}：``滴——预约核验通过。''（未预约 $\to$ 违规``翻越闸机''；进出留记录）
  \item \textbf{走廊}（幻觉事件）：``办公室的灯在闪，虽然里面没有人。'' SAN 判定：败 SAN $-10$。
  \item \textbf{办公室 · 档案柜}（钥匙）：``档案柜。今日剩余档案：1 份。'' $\to$ 指定 2 人，获知是否同一阵营（每日全图 1 份；进办公室对同社区者可见）。
  \item \textbf{社区部 · 预约簿}（钥匙）：``昨日预约登记簿。'' $\to$ 昨日预约与出入时间线。
  \item \textbf{策创部 · 修订台}（钥匙）：``校规修订台。橡皮还有最后一块。'' $\to$ 抹除 1 条随机规则（每人每局 1 次；理事长无提示）。
  \item \textbf{宣传部 · 校园墙后台}（钥匙）：``管理员模式。发布类型：破译 / 表彰 / 担保。'' $\to$ 匿名发帖 1 条/日。
  \item \textbf{咖啡屋 · 特调}（咖啡券）：``今日特调：失忆拿铁。'' $\to$ 50 HP 或 30 SAN $+5$ 饥饿。
  \item \textbf{理事长 · 假装交互}：界面与干事一致，底部多一行灰色小字``（仅供查看）''。$\to$ 查看物品清单与价值，不可拿；不耗次数。
\end{itemize}

\subsection{通用交互}
\begin{itemize}
  \item \textbf{攻击}（持杠铃杆）：同房间玩家列表浮出红框。确认 $\to$ 目标眩晕 2 min；干事额外 $-30/-30$；教学楼/图书馆内 $=$ F4 违规 + 没收。
  \item \textbf{机器查阅}（持查阅卡）：``今日记录：14:03 购买 $\times 1$；14:17 闸机 $\times 1$……'' $\to$ 1 AP + SAN 3，时间 + 次数，无人名。
  \item \textbf{私聊转账}：聊天窗口内输入金额，对方确认后到账。
  \item \textbf{检举}（理事长操作员）：目标列表 $\to$ 点选名字 $\to$ 确认。界面显示：剩余弹药 N、本日失败 M/2、预估伤害。
\end{itemize}

\clearpage
\section{挑战类题目库}
\label{app:quiz}
挑战类教室统一规则：进入后不限时答题；答对得奖励，答错 SAN $-10$；每房间每日限答 1 次；局势题以作答时刻的系统记录为准。题库可持续扩充，注意\textbf{题目与奖励不绑定固定房间}，每局随机分派。

\begin{longtable}{p{3.2cm}p{7.2cm}p{2.2cm}p{2.6cm}}
\toprule
类型 & 题目示例 & 答案 & 奖励 \\
\midrule
\endhead
脑筋急转弯 & 什么东西越洗越脏？ & 水 & 5 币 \\
脑筋急转弯 & 什么东西打破了才能吃？ & 鸡蛋 & 零食 $\times 2$ \\
脑筋急转弯 & 1+1 在什么情况下等于 3？ & 算错的时候 & 咖啡券 \\
观察题 & 此刻，这间教室里有几个人？ & 1（你自己） & 咖啡券 \\
高等数学 & $\int_0^1 x\,dx=?$ & $1/2$ & 10 币 \\
高等数学 & $e^{i\pi}+1=?$ & $0$ & 查阅卡 \\
高等数学 & $\lim_{x\to0}\frac{\sin x}{x}=?$ & $1$ & 10 币 \\
规则题 & 规则板上现在有几条随机规矩？ & 5（按当前板面） & 校规便签 \\
规则题 & 固定规矩 F4 的即时惩罚是什么？ & 没收 + HP $-10$ & 5 币 \\
局势题 & 昨天中午的投票平票了吗？ & 以系统记录为准 & 10 币 \\
局势题 & 今天至今通报过几次违规？ & 以系统记录为准 & 10 币 \\
局势题 & 你认为理事会里有几个人？ & 实际理事长存活数 & 10 币 + SAN $-5$ \\
自省题 & 你今天违反了几条规矩？ & 如实作答（填 0 且确为 0） & 护身符 \\
\bottomrule
\end{longtable}

\end{document}
